% Part: first-order-logic
% Chapter: lindstrom
% Section: introduction

\documentclass[../../../include/open-logic-section]{subfiles}

\begin{document}

\olfileid{mod}{lin}{int}
\section{پېژندګلو}

په دې څپرکي کښې موخه دا ده چې د لومړۍ درجې منطق په اړه د ليندستروم
ځانګړنه ثابته کړو: د هغو منطقونو په ډله کښې چې ځينې نور شرطونه هم
پوره کوي، دا هغه اعظمي منطق دے چې د متناهي شاهد او د لوېنهايم--سکولم
ښکته قضیې پکښې صدق کوي
(\olref[fol][com][com]{thm:compactness} او
\olref[fol][com][dls]{thm:downward-ls})۔ لومړے بايد د منطقونو هغه عمومي
ټولګی لا په عمومي ډول ځانګړے کړو چې قضيه پرې عملي کېږي۔ بحث به مو
تر \emph{اړيکيزو} ژبو محدود وي، يعنې هغو ژبو ته چې يوازې
!!{predicate}s او فردي ثابتونه لري، خو هېڅ !!{function}s نۀ لري۔

\end{document}
